% =============================================================================
%  Section: Problem Statement
% =============================================================================
\section{Problem Statement}

\subsection{Context}

An aircraft maintenance hangar contains a fixed set of \emph{positions}
$\mathcal{P} = \{p_1, \ldots, p_P\}$ where aircraft can be parked and serviced.
Each position can hold at most one aircraft at a time.
A set of aircraft $\mathcal{R} = \{r_1, \ldots, r_R\}$ must be scheduled, each
requiring a sequence of maintenance jobs carried out at a single position for its
entire stay.

The hangar has a physical access topology: some positions can only be reached by
driving through other positions.  This creates \emph{blocking arcs}: if aircraft
$r$ occupies a \emph{front} position $p$ while aircraft $r'$ needs to enter or
exit a \emph{rear} position $p'$, a ground movement is needed, generating delay
and logistical cost.

\subsection{Input Data}

\paragraph{Aircraft.}
Each aircraft $r \in \mathcal{R}$ is characterised by:
\begin{itemize}[noitemsep]
    \item $e_r \geq 0$: earliest time at which $r$ can enter the hangar.
    \item $T_r > 0$: target finish time (when the customer expects the aircraft
          back in service).  Finishing after $T_r$ incurs a delay penalty.
    \item A client identifier (for multi-client instances).
\end{itemize}

\paragraph{Jobs.}
Each aircraft $r$ has an ordered sequence of maintenance jobs
$\mathcal{J}^r = (j_1^r, j_2^r, \ldots)$ that must be executed
\emph{sequentially} at the assigned position:
\begin{itemize}[noitemsep]
    \item $d_j > 0$: processing duration of job $j$.
    \item Precedences: $j$ must finish before $j'$ starts whenever $(j,j')
          \in \mathcal{E}$.  Within a single aircraft the jobs form a chain;
          cross-aircraft precedences are also possible.
\end{itemize}
The \emph{total duration} of aircraft $r$ is $D_r = \sum_{j \in \mathcal{J}^r} d_j$.

\paragraph{Hangar topology.}
The blocking structure is a directed acyclic graph (DAG):
\[
    \mathcal{B} \subseteq \mathcal{P} \times \mathcal{P},
    \quad (p, p') \in \mathcal{B} \;\Rightarrow\;
    \text{an aircraft in } p \text{ blocks access to } p'.
\]
A position's \emph{blocking load} is
$\ell(p) = |\{p' : p \text{ can reach } p' \text{ transitively in } \mathcal{B}\}|$.

\paragraph{Separation.}
A minimum time gap $\delta \geq 0$ must be observed between consecutive aircraft
occupying the same position: if aircraft $r$ finishes at $f_r$ in position $p$,
the next aircraft $r'$ assigned to $p$ must start no earlier than
$f_r + \delta$.

\subsection{Decisions}

\begin{enumerate}[noitemsep]
    \item \textbf{Assignment}: which position $p \in \mathcal{P}$ is assigned to
          aircraft $r$.  Each aircraft occupies exactly one position.
    \item \textbf{Scheduling}: start time $s_j$ for each job $j$, respecting
          precedences, earliest-start, and non-overlap within each position.
\end{enumerate}

\subsection{Objective Function}

The objective is a weighted sum of three KPIs:

\begin{equation}
    \min \quad
    w_M \cdot M
    + w_D \cdot \sum_{r \in \mathcal{R}} \Delta_r
    + w_V \cdot V
    \label{eq:objective}
\end{equation}

\noindent where:

\begin{center}
\begin{tabular}{lll}
\toprule
Symbol & Definition & Default weight \\
\midrule
$M$ & Makespan: $\max_{r} f_r$ (time until the last aircraft leaves) & $w_M = 0.1$ \\
$\Delta_r$ & Individual delay: $\max(0,\; f_r - T_r)$ & $w_D = 1.0$ \\
$V$ & Total blocking movements count & $w_V = 10$ \\
\bottomrule
\end{tabular}
\end{center}

The weights above correspond to the experimental configuration used in this
project.  The MILP model uses higher values ($w_M=10$, $w_D=100$, $w_V=1$)
to avoid numerical issues; all heuristics use the same weights as the MILP when
called from the experiment runner.

\subsection{Blocking Movements}

A \emph{blocking event} occurs when aircraft $r$ arrives at or departs from
rear position $p'$ while aircraft $r_b$ is occupying a front position $p$ such
that $(p,p') \in \mathcal{B}$ and their time windows overlap.

Formally, let $s_r$, $f_r$ denote the start and finish of aircraft $r$:

\begin{itemize}[noitemsep]
    \item \textbf{Entry block}: $r'$ enters $p'$ while $r_b \in p$:
          $s_{r'} \in (s_{r_b},\; f_{r_b} - \delta)$.
    \item \textbf{Exit block}: $r'$ exits $p'$ while $r_b \in p$:
          $f_{r'} \in (s_{r_b},\; f_{r_b} - \delta)$.
\end{itemize}

Each detected blocking event contributes 2 to $V$ (one movement in, one out).

\subsection{Feasibility}

A solution is \emph{feasible} if:
\begin{enumerate}[noitemsep]
    \item Every aircraft is assigned to exactly one position.
    \item Jobs of an aircraft are scheduled at the assigned position, in order,
          without gaps within the position chain.
    \item No two aircraft overlap at the same position (minimum separation $\delta$
          between consecutive aircraft).
    \item Each aircraft's first job starts no earlier than $e_r$.
    \item All job precedences $(j, j') \in \mathcal{E}$ are satisfied.
\end{enumerate}
